%  08_ai_usage.tex  –  Use of Artificial Intelligence

\chapter*{Use of Artificial Intelligence}
\addcontentsline{toc}{chapter}{Use of Artificial Intelligence}
\label{ch:ai-usage}

Throughout the lifecycle of this academic project, artificial intelligence (AI) was integrated
as a support assistant to enhance productivity across both the technical implementation and the
documentation phases. The tool utilized was the Large Language Model
Claude~\cite{claudeai}.

In the software development phase, AI assistance was used for code improvement, backend logic
and user interface design. Specifically, the tool supported the creation of the edge computing
script \texttt{predict.py}, which is deployed via \textit{cron} on the Raspberry Pi 4 to
compute daily 7-day air quality forecasts, as well as the complete machine learning training
pipeline executed within the Google Colab notebook
\texttt{AirQuality\_ML\_Training.ipynb}. On the server side, the backend prediction and
administrative messaging endpoints were integrated into \texttt{web\_db\_connection.py} with
AI support. Furthermore \texttt{predictions.html} was developed using the tool. It also helped
with the application's cascading stylesheet (\texttt{style.css}), which accelerated the
implementation of layout architecture, responsive design, and component-level styling.

Another use for the AI assistant was as a diagnostic asset for root-cause analysis during
system debugging. It proved instrumental in identifying and resolving critical bugs or
irregularities that occurred during development.

It also helped navigating different AWS features and how to set them up correctly. Furthermore,
AI was utilized for suggestions on how to structure a research paper.

AI was also used to speed up the process of creating the Android mobile application. Its help
was needed in places where something new appeared or some technical details had to be resolved,
such as: creation of the HTTP API service layer for connection between the app and Flask server
and configuration of plain HTTP traffic support through \texttt{network\_security\_config.xml}.
It also assisted with \texttt{android:usesCleartextTraffic} attribute of manifest file and the
implementation of bottom navigation bar with tab states which are saved based on
\texttt{IndexedStack}.

In conclusion, AI was used with the role of assistant. It accelerated development, helped
identify bugs and assisted with the documentation. All AI-generated code was reviewed manually
and tested against the actual hardware and server environment before being considered complete.
